\documentclass{article}
\usepackage[T1]{fontenc}
\usepackage{lmodern}
\usepackage{graphicx}
\usepackage{amsmath,amssymb}
\usepackage[a4paper,margin=2cm]{geometry}
\usepackage[absolute,overlay]{textpos}
\setlength{\TPHorizModule}{1mm}
\setlength{\TPVertModule}{1mm}
\usepackage[indonesian]{babel}
\usepackage{hyperref}
\setlength{\emergencystretch}{2em}
% unit: Halaman 17

\begin{document}

% === Halaman 17 (llm) ===

\section*{Galois Field GF($2^m$)}

\begin{itemize}
    \item Dari GF($p^m$), jika $p = 2$, maka GF($2^m$) disebut juga medan berhingga biner.
    \item GF($2^m$) atau $\text{F}_2^m$ adalah ruang vektor berdimensi $m$ pada GF(2). Setiap elemen di dalam GF($2^m$) adalah integer dalam representasi biner sepanjang maksimal $m$ bit.
    Contoh: GF($2^4$), GF($2^8$), dst
    \item String biner $b_{m-1}b_{m-2} \dots b_1b_0$, $b_i \in \{0,1\}$, dapat dinyatakan sebagai polinom
    \[ b_{m-1}x^{m-1} + b_{m-2}x^{m-2} + \dots + b_1x + b_0 \]
    \item Jadi, setiap $a \in \text{GF}(2^m)$ dapat dinyatakan sebagai
    \[ b_{m-1}x^{m-1} + b_{m-2}x^{m-2} + \dots + b_1x + b_0 \]
    \item Contoh: 1101 di dalam GF($2^4$) dapat dinyatakan sebagai polinom $x^3 + x^2 + 1$
    
    0101010111 di dalam GF($2^8$) dapat dinyatakan sebagai polinom $x^6 + x^4 + x^2 + x + 1$
\end{itemize}

\end{document}
